% 三面板调用同一个街景函数，位置与轮廓严格共用。
\begingroup%
% 第1—3章局部矢量图元：单色黄、双色蓝红、三色蓝红黄。
\providecommand{\BasicsFiveSetup}{%
 \colorlet{B5Blue}{FigureBlue}\colorlet{B5Coral}{FigureRed}%
 \colorlet{B5Yellow}{FigureYellow}\definecolor{B5Gray}{HTML}{4C4D4F}%
 \colorlet{B5Ice}{FigureBlueFill}\colorlet{B5Rose}{FigureRedFill}%
 \colorlet{B5Cream}{FigureYellowFill}%
 \tikzset{b5 picture/.style={x=1mm,y=1mm,font=\figurefont\mdseries\fontsize{8.5}{10.5}\selectfont,text=B5Gray,line width=1pt},%
 b5 node/.style={draw=B5Gray,fill=white,line width=1pt,rounded corners=1.5mm,inner sep=3pt,font=\figurefont\mdseries\fontsize{8.5}{10.5}\selectfont,text=B5Gray,align=center},%
 b5 flow/.style={draw=B5Gray,line width=1.1pt,-{Stealth[length=2.2mm,width=1.4mm]}},%
 b5 link/.style={draw=B5Gray,line width=.7pt},%
 b5 feedback/.style={b5 flow,draw=B5Coral,dash pattern=on 2.5mm off 1.5mm}}}%
% 3—3—1网络仅为图标，不指定真实架构；各层用实心蓝、红、黄，连线中性灰。
\providecommand{\BasicsFiveNet}[3]{\begin{scope}[shift={(#1,#2)},scale=#3]%
 \foreach \a in {-3,0,3}{\foreach \b in {-3,0,3}{\draw[b5 link] (0,\a)--(5,\b);}}%
 \foreach \a in {-3,0,3}{\draw[b5 link] (5,\a)--(10,0);}%
 \foreach \a in {-3,0,3}{%
  \draw[draw=B5Gray,line width=.8pt,fill=B5Blue] (0,\a) circle (1.1);%
  \draw[draw=B5Gray,line width=.8pt,fill=B5Coral] (5,\a) circle (1.1);}%
 \draw[draw=B5Gray,line width=.8pt,fill=B5Yellow] (10,0) circle (1.1);%
\end{scope}}%
% 鸟与猫保留教学对象，不生成或模拟真实数据样本。
\providecommand{\BasicsFiveBird}[4]{\begin{scope}[shift={(#1,#2)},scale=#3]%
 \draw[draw=B5Gray,fill=#4,line width=.8pt,line join=round]%
  (-6,-1)--(-10,-2)--(-6,2).. controls (-3,3) and (-2,6) .. (1,6)%
  .. controls (4,6) and (5,4) .. (4,2).. controls (4,-1) and (-1,-3) .. (-6,-1);%
 \draw[draw=B5Gray,fill=white,line width=.7pt] (4,3)--(6,3.8)--(4,4.1)--cycle;%
 \fill[B5Gray] (2.7,4.3) circle (.35);%
 \draw[b5 link] (-5,1).. controls (-1,4) and (1,1) .. (-3,-.5);%
 \draw[b5 link] (-2,-2)--(-2,-4)--(-.7,-4);\draw[b5 link] (0,-1.5)--(0,-4)--(1.3,-4);%
\end{scope}}%
\providecommand{\BasicsFiveCat}[3]{\begin{scope}[shift={(#1,#2)},scale=#3]%
 \draw[draw=B5Gray,fill=white,line width=.8pt] (0,-1) ellipse (3.3 and 4.3);%
 \draw[draw=B5Gray,fill=white,line width=.8pt,line join=round]%
  (-3,3)--(-3,7)--(-1.2,5.8).. controls (0,6.5) and (1,6.5) .. (2,5.8)%
  --(3.3,7)--(3.3,3).. controls (2,1.3) and (-1.5,1.3) .. (-3,3);%
 \fill[B5Gray] (-1.3,4) circle (.3);\fill[B5Gray] (1.5,4) circle (.3);%
 \draw[b5 link] (0,3.4)--(0,2.7);\draw[b5 link] (-2,-2)--(-2,-5);\draw[b5 link] (2,-2)--(2,-5);%
 \draw[draw=B5Gray,line width=.8pt] (-2.5,-3).. controls (-8,-5) and (-8,-.5) .. (-5,-1);%
\end{scope}}%
\providecommand{\BasicsFiveRobot}[3]{\begin{scope}[shift={(#1,#2)},scale=#3]%
 \fill[B5Gray,rounded corners=.4mm] (-2.7,-1.6) rectangle (-1.7,1.6);%
 \fill[B5Gray,rounded corners=.4mm] (1.7,-1.6) rectangle (2.7,1.6);%
 \draw[draw=B5Gray,fill=white,line width=1pt,rounded corners=.7mm] (-1.8,-2.3) rectangle (1.8,2.3);%
 \draw[draw=B5Gray,fill=B5Yellow,line width=.8pt] (0,.6) circle (.65);%
\end{scope}}%
\BasicsFiveSetup%
\definecolor{B5Building}{HTML}{E9E9E9}%
\colorlet{B5Sign}{FigureYellow}%
\newcommand{\BStreet}[2]{\begin{scope}[shift={(#1,0)}]%
 \ifnum#2=2%
  \def\CarFill{B5Blue}\def\WindowFill{B5Blue}\def\WheelFill{B5Blue}%
  \def\BodyFill{B5Coral}\def\SignFill{B5Sign}\def\RoadFill{B5Cream}%
 \else%
  \def\CarFill{white}\def\WindowFill{B5Ice}\def\WheelFill{B5Gray}%
  \def\BodyFill{white}\def\SignFill{white}\def\RoadFill{B5Ice}%
 \fi%
 \fill[B5Ice] (0,0) rectangle (52,50);%
 \fill[white] (0,0) rectangle (52,31);%
 \draw[draw=B5Gray,fill=\RoadFill,line width=.8pt]%
  (0,0)--(0,7)--(19,31)--(33,31)--(52,7)--(52,0)--cycle;%
 \draw[draw=B5Gray,fill=B5Building,line width=.8pt]%
  (0,32)--(0,44)--(7,44)--(11,40)--(11,32)--cycle;%
 \draw[b5 link] (7,44)--(7,32);%
 \draw[draw=B5Gray,fill=B5Building,line width=.8pt]%
  (41,32)--(41,40)--(45,44)--(52,44)--(52,32)--cycle;%
 \draw[b5 link] (45,44)--(45,32);%
 \foreach \xx in {6,47}{%
  \draw[draw=B5Gray,fill=white,line width=.8pt]%
   (\xx-5,29).. controls (\xx-6,31) and (\xx-4,32) .. (\xx-3,32)%
   .. controls (\xx-3,36) and (\xx+1,36) .. (\xx+2,33)%
   .. controls (\xx+4,33) and (\xx+5,31) .. (\xx+5,29)--cycle;}%
 \ifnum#2<2%
  \draw[white,line width=1.1pt,dash pattern=on 2mm off 1.5mm] (26,29)--(26,20);%
 \fi%
 \draw[draw=B5Gray,fill=\SignFill,line width=.8pt] (8.6,5) rectangle (9.4,19);%
 \draw[draw=B5Gray,fill=\SignFill,line width=.8pt] (9,21) circle (2.6);%
 \draw[draw=B5Gray,fill=\WheelFill,line width=.8pt,rounded corners=.5mm] (19.5,4) rectangle (22,8);%
 \draw[draw=B5Gray,fill=\WheelFill,line width=.8pt,rounded corners=.5mm] (30,4) rectangle (32.5,8);%
 \draw[draw=B5Gray,fill=\CarFill,line width=.8pt,rounded corners=.6mm]%
  (18,7)--(18,13)--(21,17)--(31,17)--(34,13)--(34,7)--cycle;%
 \draw[draw=B5Gray,fill=\WindowFill,line width=.7pt] (20,12)--(21.5,16)--(30.5,16)--(32,12)--cycle;%
 \foreach \xx in {19,31}{\draw[draw=B5Gray,fill=\CarFill,line width=.7pt] (\xx,9.5) rectangle (\xx+2,10.5);}%
 \draw[draw=B5Gray,fill=\CarFill,line width=.7pt] (24.5,8) rectangle (27.5,9);%
 \draw[draw=B5Gray,fill=\BodyFill,line width=.8pt] (44,18.5) circle (1.4);%
 \draw[draw=B5Gray,fill=\BodyFill,line width=.8pt,line join=round]%
  (43,16).. controls (44,16.5) and (45,15.5) .. (46,13)%
  --(47,10.5)--(46,10)--(44.8,12.2)--(45,8.5)--(47,4.5)--(46.2,4)%
  --(43.8,8)--(42.8,4)--(41.6,4)--(42.8,9.5)--(43,12.2)%
  --(41,10)--(40.3,10.5)--(42,14)--cycle;%
 \ifnum#2=1%
  \draw[draw=B5Coral,line width=1pt] (6,4) rectangle (12,24);%
  \draw[draw=B5Coral,line width=1pt] (17,3.5) rectangle (35,18);%
  \draw[draw=B5Coral,line width=1pt] (40,3.5) rectangle (48,20);%
  \node at (9,27) {标志};\node at (26,21) {车辆};\node at (44,23) {行人};%
 \fi%
 \ifnum#2=2%
  \node at (26,44) {天空};\node at (3.5,39) {建筑};\node at (48.5,39) {建筑};%
  \node at (6,31) {植被};\node at (47,31) {植被};\node at (26,28) {道路};%
  \node at (26,21) {车辆};\draw[b5 link] (26,19.5)--(26,17.5);%
  \node at (45,25) {行人};\draw[b5 link] (45,23.5)--(44.5,20.5);%
  \node[anchor=west] at (1,26) {标志};\draw[b5 link] (6,24.5)--(7,23);%
 \fi%
 \draw[draw=B5Gray,line width=1pt] (0,0) rectangle (52,50);%
\end{scope}}%
\begin{tikzpicture}[b5 picture]%
 \path[use as bounding box] (0,-1) rectangle (168,59);%
 \node[anchor=west] at (0,56) {输入街景};%
 \node[anchor=west] at (58,56) {目标检测};%
 \node[anchor=west] at (116,56) {语义分割};%
 \BStreet{0}{0}\BStreet{58}{1}\BStreet{116}{2}%
\end{tikzpicture}%
\endgroup%
